\begin{vidu}
Một vật dao động điều hoà có phương trình gia tốc: $a=12\pi^2\cos(2\pi t+\dfrac{\pi}{2})\ cm/s^2$.
\begin{itemize} 
\item[1.] Tần số góc của dao động là\dotfill
\item[2.] Chu kì của dao động là\dotfill
\item[3.] Tần số của dao động là\dotfill
\item[4.] Biên độ của dao động là\dotfill
\item[5.] Phương trình dao động của vật có dạng\dotfill
\item[6.] Phương trình vận tốc của vật có dạng\dotfill
\end{itemize}
\loigiai{
\begin{itemize}
\item[a)] $\omega=2\pi\ rad/s;\ A=3\ cm;\ \varphi=-\dfrac{\pi}{2};\ T=1\ s;\ f=1\ Hz$. 
\item[b)] $x=3\cos(2\pi t-\dfrac{\pi}{2});\ v=6\pi\cos(2\pi t)\ cm/s$. 

\end{itemize}
}
\end{vidu} %\dotlineEX{5}

%\loai{Tính vận tốc cực đại và gia tốc cực đại}
\begin{vidu}
(QG - 2021) Một chất điểm dao động điều hòa với phương trình $x=4\cos5t\ (cm)$ (t tính bằng s). Tốc độ của chất điểm khi đi qua vị trí cân bằng là
\choice
{50 cm/s}
{\True 20 cm/s}
{100 cm/s}
{80 cm/s}
\loigiai{
\begin{itemize}
\item $v_{\max}=\omega A=5.4=20\ cm/s$.
\end{itemize}

}
\end{vidu} 
%\dotlineEX{2}
\begin{vidu}%[3P1-2.1B][Loai1]
Một vật nhỏ dao động điều hòa trên quỹ đạo dài 20 cm với tần số góc là 6 rad/s. Gia tốc cực đại của vật có giá trị là
\choice
{7,2 $m/s^2$}
{0,72 $m/s^2$}
{\True 3,6 $m/s^2$}
{0,36 $m/s^2$}
\loigiai{
Ta có: $A = \dfrac{L}{2} = 10\ cm \Rightarrow a_{\max} = \omega^2A = 360\ cm/s^2 = 3,6\ m/s^2$.
%\dotlineEX{2}
}
\end{vidu}
%\dotlineEX{2}
%\loai{Từ vận tốc cực đại hoặc gia tốc cực đại tìm các đại lượng}

\begin{vidu}
Một chất điểm dao động điều hòa với phương trình $x=4 \cos (\pi t+\dfrac{\pi}{3})\ \mathrm{cm}$.
\begin{itemize}
\item[1.] Biên độ dao động của chất điểm là\dotfill
\item[2.] Quỹ đạo dao động của chất điểm là\dotfill
\item[3.] Pha dao động của chất điểm tại thời điểm $t$ là\dotfill
\item[4.] Pha ban đầu của chất điểm là\dotfill
\item[5.] Pha dao động của chất điểm tại $t=\dfrac{1}{3}\ s$ là\dotfill
\item[6.] Tần số góc của chất điểm là\dotfill
\item[7.] Tần số dao động của chất điểm là\dotfill
\item[8.] Chu kì dao động của chất điểm là\dotfill
\item[9.] Số dao động toàn phần chất điểm thực hiện được trong $10 \mathrm{~s}$ là\dotfill
\item[10.] Li độ cực đại của chất điểm là\dotfill
\item[11.] Li độ cực tiểu của chất điểm là\dotfill
\item[12.] Tốc độ cực đại của chất điểm là\dotfill
\item[13.] Tốc độ cực tiểu của chất điểm là\dotfill
\item[14.] Vận tốc cực đại của chất điểm là\dotfill
\item[15.] Vận tốc cực tiểu$ $ của chất điểm là\dotfill
\item[16.] Gia tốc của chất điểm có giá trị cực đại là\dotfill
\item[17.] Gia tốc của chất điểm có giá trị cực tiểu là\dotfill
\item[18.] Gia tốc của chất điểm có độ lớn cực đại là\dotfill
\item[19.] Gia tốc của chất điểm có độ lớn cực tiểu là\dotfill
\item[20.] Phương trình vận tốc của chất điểm là\dotfill
\item[21.] Phương trình gia tốc của chất điểm là\dotfill
\item[22.] Vận tốc của chất điểm tại $t=\dfrac{10}{3} \mathrm{~s}$ là\dotfill
\item[23.] Gia tốc của chất điểm tại $t=\dfrac{20}{3} \mathrm{~s}$ là\dotfill
\end{itemize}
\loigiai{
\begin{center} 
\begin{tabular}{|l|l|l|}
\hline 1. $A=4\ cm$ & 2. $L=2 A=8\ cm$ & 3. $\varphi=\pi t+\dfrac{\pi}{3}\ rad$ \\
\hline 4. $\varphi=\dfrac{\pi}{3}\ rad$ & 5. $\varphi=\dfrac{\pi}{3}+\dfrac{\pi}{3}=\dfrac{2\pi}{3}\ rad$ & 6. $\omega=\pi\ rad/s$ \\
\hline 7. $f=\dfrac{\omega}{2\pi}=0,5\ Hz$ & 8. $T=\dfrac{1}{f}=2\ s$ & 9. $10\ s=5 T \rightarrow 5\ dd$ \\
\hline 10. $x_{\max }=A=4\ cm$ & 11. $x_{\min }=-A=-4\ cm$ & 12. $|v|_{\max }=\omega A=4 \pi\ cm/s$ \\
\hline $13 .|v|_{\min }=0\ cm/s$ & 14. $v_{\max }=\omega A=4 \ cm/s$ & 15. $v_{\min }=-\omega A=-4 \pi\ cm/s$ \\
\hline $16 .\ a_{\max }=\omega^2 A=4 \pi^2\ cm/s^2$ & 17. $a_{\min }=-\omega^2 A=-4 \pi^2\ cm/s^2$ & $18 . x=4 \cos (\dfrac{2 \pi}{3}+\dfrac{\pi}{3})=-4$ \\
\hline $19 .\ v=4 \pi \cos (\pi t+\dfrac{5 \pi}{6})$ & $20 . a=4 \pi^2 \cos (\pi t-2 \pi / 3)$ & $21 . v=2 \pi \sqrt{3}\ cm/s$ \\
\hline 22. $a=4 \pi^2\ cm/s^2$ & 
 & \\
 \hline
\end{tabular}
\end{center} 
}
\end{vidu}
\begin{vidu}%[3P1-2.1B][Loai1]
Một vật nhỏ dao động điều hòa với biên độ 10 cm và vận tốc có độ lớn cực đại là $10\pi$  cm/s. Tần số dao động bằng
\choice
{$\pi$  Hz}
{\True 0,5 Hz}
{1 Hz}
{2 Hz}
\loigiai{
Ta có: $v_{\max} = \omega A$ $\Rightarrow$ $\omega  = \dfrac{v_{\max}}{A} = \pi  \ rad/s \Rightarrow f = 0,5$ Hz.

}
\end{vidu} 
%\dotlineEX{2}
%\begin{vidu}
%Một máy cơ khi hoạt động sẽ tạo ra những dao động được xem gần đúng là dao động điều hoà với phương trình dao động có dạng: $x=0,1\cos(180\pi t)\ mm$. 
%\begin{itemize} 
%\item[1.] Biên độ của dao động là ............................... 
%\item[2.] Chu kì của dao động là ...............................
%\item[3.] Tần số của dao động là ...............................
%\item[4.] Tần số góc của dao động là ...............................
%\item[5.] Vận tốc cực đại của dao động là ...............................
%\item[6.] Gia tốc cực đại của dao động là ...............................
%\item[7.] Phương trình vận tốc của dao động là ...............................
%\item[8.] Phương trình gia tốc của dao động của dao động là ...............................
%\item[9.]  Vận tốc tại thời điểm $t=\dfrac{1}{45}\ s$ của dao động là ...............................
%\item[10.]  Gia tốc tại thời điểm $t=\dfrac{1}{45}\ s$ của dao động là ...............................
%\end{itemize}
%\loigiai{
%\begin{itemize}
%\item[a)] Biên độ: $A=0,1\ mm$; Tần số góc: $\omega=180\pi\ rad/s$; Tần số: $f=90\ Hz$; Chu kì: $\dfrac{1}{90}\ s$.
%\item[b)] $v_{\max}=\omega A=18\pi\ mm/s;\ a_{\max}=\omega^2 A=3240\pi^2\ mm/s^2\approx 32\ m/s^2$. 
%\item[c)] Phương trình vận tốc: $v=-18\pi\sin(180\pi t)\ mm/s$\\
%Phương trình gia tốc: $a=-3240\pi^2\cos(180\pi t)\ mm/s^2$.
%\item[d)] Tại $t=\dfrac{1}{45}\ s=\dfrac{T}{2}$ vật đang ở vị trí biên dương nên $v=0$ và $a=a_{\min}=-3240\pi^2\ mm/s^2$.
%
%\end{itemize}
%}
%%\dotlineEX{6}
%\end{vidu} %\dotlineEX{2}

%\loai{Từ vận tốc cực đại và gia tốc cực đại tìm các đại lượng}
\begin{vidu}%[3P1-2.1K][Loai1]
Một vật dao động điều hòa dọc theo trục Ox.Vận tốc cực đại của vật là $v_{\max} = 4\pi$  cm/s và gia tốc cực đại $a_{\max} = 8\pi ^2$ $ cm/s^2$. Quỹ đạo dao động có độ dài là
\choice
{8 cm}
{2 cm}
{16 cm}
{\True 4 cm}
\loigiai{
Ta có: $\begin{cases}v_{\max} = 4\pi \ cm/s = \omega A\\ \ a_{\max} = 8\pi ^2 = \omega^2A\end{cases} \Rightarrow \omega  = 2\pi \Rightarrow A = 2\ cm \Rightarrow L = 2A = 4$ cm.
}
\end{vidu} %\dotlineEX{2}
%\loai{Tìm vận tốc và gia tốc tại thời điểm t}
\begin{vidu}%[3P1-2.1K][Loai1] 
Một chất điểm dao động điều hòa theo phương trình $x = 2\cos\left(\pi t - \dfrac{2\pi}{3}\right)$ cm. Tại thời điểm $t = 7,5$ s chất điểm đang chuyển động với
\choice
{tốc độ $2\pi$ cm/s theo chiều dương}
{tốc độ $\pi$ cm/s theo chiều dương}
{\True tốc độ $\pi$ cm/s theo chiều âm}
{tốc độ $2\pi$ cm/s theo chiều âm}
\loigiai{
\begin{itemize}
\item $x = 2\cos\left(\pi t-\dfrac{2\pi}{3}\right)$ cm
$\Rightarrow v = -2\pi\sin\left(\pi t-\dfrac{2\pi}{3}\right)$ cm/s.
\item Tại thời điểm $t = 7,5$ s, $v = -\pi\ cm/s < 0\Rightarrow$ chất điểm đang chuyển động với tốc độ $\pi$ cm/s theo chiều âm.
\end{itemize}}
\end{vidu} %\dotlineEX{2}

%\loai{Tổng hợp}
\begin{vidu}%[3P1-2.1K][Loai1]
Một chất điểm dao động điều hòa với phương trình gia tốc $a = 100\cos(5t + \dfrac{\pi}{3})$ (a tính bằng $ cm/s^2$, t tính bằng s). Phát biểu nào sau đây đúng?
\choice
{Tốc độ cực đại của chất điểm là 10 cm/s}
{Gia tốc của chất điểm có độ lớn cực đại là 500 $ cm/s^2$}
{Tần số của dao động là 5 Hz}
{\True Biên độ dao động là 4 cm}
\loigiai{
\begin{itemize}
\item Ta có: $a = 100\cos(5t + \dfrac{\pi}{3})$ $\Rightarrow$ $a_{\max} = \omega^2A = 100\ cm/s^2;\ \omega  = 5\ rad/s \Rightarrow A = 4$ cm.
 
\end{itemize}
}
\end{vidu} %\dotlineEX{4}